% !TeX root = ../papers/nc-p3.tex
\section{片側だけの共通切断と有限な反対階数九の判定}
\label{sec:finite-rank-nine-all-marks}
\begin{lemma}\label{lem:quartic-two-marks-incidence}
全中央支持4の有限窓が、二つの非伸長型平面局所点 $p,q$ に支持されるとする。
各点で少なくとも一方に共通四空間があり、その階数一切断が長さ4であるなら、
少なくとも片側の反対対称性が成立する。二点が選ぶ側は異なってよい。
\end{lemma}
\begin{proof}
対称性がないとし、まず $p$ の共通切断の核側の一座標で三次持上げを取る。
定理\ref{thm:finite-artin-kernel-symmetry}の外積の支持 $Z$ は長さ3または4である。
共線の場合は一つの巡回行列による純関係で排除される。
長さ4の曲線状因子の積は系\ref{cor:single-support-quartic-slice}の三重平面または
印で重複度3の行列式による矛盾で除かれる。この計算には反対切断階数も他側の一致も使わない。
残るのは曲線状因子の積の長さ3または局所 $(1,2,1)$ 型である。
$Z$ は $p$ を含む。実際、含まなければ定理\ref{thm:finite-stretched-planar-marks}の
小行列式による $\mathfrak m_p^2 f\subset(\eta)$ の評価で
隣接語 $f=H\chi(a)$ の印での支持長は高々2、外では $(\eta)$ に入る。
長さ3なら全長6の下界に反し、長さ4で全長6に達すれば $f$ の支持が
偶座標で全空間を張る $Z=\Sigma$ を含み、長さ6の平面条件に反する。

三次持上げの一次因子 $H$ は $p$ で非零である。実際印でのsocle分離式は
$H$ と印からの二つの射影一次式の積で、後の二式は印で消える。
行列式の印での重複度はちょうど2である（系\ref{cor:single-support-quartic-slice}の二次初期式の計算）。
一方 $H$ は点環全体の単元ではない。単元なら、逆数分子の零でない二次座標が
$(\eta)$ に入り、その支持環 $Z$ には隣の偶座標が全射する。隣り合う座標の二次空間どうしの
完全跡直交からこの二次式は零になり、二次制限の単射性に反する。
よって $H(q)=0$。$Z$ が局所なら印での分離式は $Hs^2$ である。
$Z$ の長さが $2+1$ に分かれれば、単純因子の分離式が $Hs^2$ となる。
定理\ref{thm:finite-split-separator-slices}よりその因子の閉点で $H$ は非零であるから、
単純因子の閉点は必ず $p$ である。従ってこの場合も $p$ の分離式が $Hs^2$ である。
核の二位置を交換した同じ議論を適用できる。これで二つの核位置のそれぞれに同じ印 $p$ の重平面分離式がある。

長さの大きい方を $q$ と取る。その長さは10以上なので、どの位置でも $H$ は
$q$ の局所環で恒等零ではない（定理\ref{thm:finite-plane-section-seven}）。
各重平面分離式が $q$ の局所環で零だから $s(q)=0$ であり、これが当該位置の
$p$ の接平面である。実際重平面模型 $E_D$ の $D$ は交代で、その $s$ 座標が
極大イデアルの二乗に入ることは命題\ref{prop:finite-two-planar-tangent-incidence}の
$s\mathfrak m\subset\mathfrak m^3$ の計算による。隣の二位置でも、同模型
$A=H vv^t+K$（$K$ は交代）の左右核から $H(v^tl)^2=H(v^tr)^2=0$ となる。
左右の基底はこの模型に従って取り、$v^tl,v^tr$ はその二位置の印の接平面の式である。
よって $q_i\in T_{p_i}X_i$ が四位置全てで成立する。この推論には反対階数8での
二位置の接平面入射の同値を使わない。

今度は $q$ の共通切断で同じ証明を行う。$q$ における各重平面分離式の単因子は
$q$ で非零なので、$p$ の接平面とは異なる。従ってこの一次式は $p$ の局所環で
恒等零とはなり得ない。実際、恒等零の平面は二次元接平面を含み、これに等しくなる。
$p$ の長さに関係なく上の核計算から $p_i\in T_{q_i}X_i$ が成立する。
両端で接する直線は切断長4以上となり系\ref{cor:finite-no-four-on-line}に反する。
\end{proof}

\begin{lemma}\label{lem:high-order-point-rank-nine}
有限窓で全中央支持が4、点スキームが一つの平面局所点に支持され、
その二つの定義式がともに四次以上なら、四位置を揃えると共通八点切断があり、
窓は補題\ref{lem:common-net-alternating-part}の分解を持つ。
四形式が非零なら反対階数は少なくとも12である。
\end{lemma}
\begin{proof}
ここでは反対位置での二つの四空間が一致することを仮定しない。
一つの文字を印で固定し、反対位置の変数 $u$ による四空間 $A(u)$ を一つずつ取る。
隣の二座標はそれぞれ $A(u)$ の核を与える。
定理\ref{thm:finite-high-order-common-octad}の一次ブロック計算は
この一行列だけを使い、印での二つの接方向では右下ブロックが零になる。
そのブロックへの射影の核はちょうど三次元である。全四空間がそこにあれば
下の行からの非零評価の階数が1以下になる。ブロック像の生成元 $D$ も同じ理由で可逆である。
印の階数一切断の接次元は少なくとも2で、正次元切断ではあり得ない。
実際定理\ref{thm:finite-marked-ruling-pencil}の母線だけとなる証明の前半は
四空間単独の左右評価階数の下界から従い、母線では全点で接次元が1である。
従って印以外の二つの接行列の上行への射影も左列への射影も独立であり、
さもなければ階数一行列の直線ができる。これで各四空間は個別に $E_D$ のブロック表示を持つ。
ここでは対称な $D$ も許す。

各辺では、残る二位置の一方を印に固定して得る三関係ずつを取る。
各組の下二式は両位置の印での評価が零の関係空間を張る。これは辺の一次部分の階数4から
二次元であり、二組で共有される。従って各辺に定理\ref{thm:finite-high-order-common-octad}の
四式 $xm=ny,\ sm=q(x,y),\ nt=r(x,y)$ がある。ここまでに同じ位置を変数とする
別の四空間との同一視は要らない。その後の三次初期式の比較、局所完全交叉の連結、
$n_i^2$ の零化イデアルによる四辺の同定はこの四式のみを使う。唯一の支持が印なので
八点切断は全辺で一致し、補題\ref{lem:common-net-alternating-part}の分解を得る。

定理\ref{thm:high-order-octad-bisymmetry}の既約二次曲面上の左右核の証明も、
$A(u)$ だけの行列式とその両隣の辺だけを使う。従って $A(u)$ の対称部分は三次元以上を張る。
これらは奇側の $p$ を消す対称形式である。奇偶を交換して同じ結論を得る。
$G$ の反対切断をその階数だけの空間と双対との完全対合に分解すると、
偶側で $p$ と任意の $u$ を入れた像と、奇側で $p$ と任意の $v$ を入れた像は直交する。
実際 $w(p,p,u,v)=0$ は辺の関係であり、四形式もこの評価では零である。
両方がそれぞれ三次元以上なので $G$ の階数は6以上である。
四形式が非零なら独立な外積ブロックの階数6を加えて下界12を得る。
\end{proof}

\begin{theorem}\label{thm:finite-rank-nine-all-marks}
四辺が有限な反対階数9の標準有限Frobenius候補が標準Koszul左右AS正則となる
必要十分条件は、中央支持3の点が存在することである。印の被約性を仮定しない。
\end{theorem}
\begin{proof}
中央支持3があれば定理\ref{thm:finite-cubic-separators}を使う。残りでは全て4である。
各印で、偶二位置を別々に印に固定した二つの四空間の和と、
奇二位置に対する同様の和は、反対切断の完全対合で相互に直交する。
これは四つの隣接点対全てでの消滅である。和の次元の合計は9以下であり、
各々4以上だから少なくとも一方は4、その側に共通切断がある。

対称性がどちらにもないとする。曲線状の印、伸長型の印、長さ6以下の印での
三重点・四点の各証明は一つの共通切断だけを使う。長さ3ではその変数側、
長さ4では核側で対称性を与える（定理\ref{thm:finite-artin-kernel-symmetry}、
\ref{thm:finite-stretched-planar-marks}、\ref{thm:finite-nodal-cuspidal-marks}の各局所計算）。
長さ2以下と正次元の場合も一つの共通切断で足りる。
従って残る点は非伸長型平面局所点で各長さ7以上、台は高々二点となる。
四次以上で始まる局所完全交叉は長さ16以上なので唯一の台となり、
補題\ref{lem:high-order-point-rank-nine}から両対称である。
それ以外の印では局所曲線の最小重複度が2または3である。
共通切断が長さ3なら、重複度3でその切断は二次初期式の不在から平方零型であり、
いずれも補題\ref{lem:finite-carrier-blowup}で対称になる。
よって各点に長さ4の共通切断がなければならない。一点支持なら
系\ref{cor:single-support-quartic-slice}、二点支持なら
補題\ref{lem:quartic-two-marks-incidence}で対称性が従い、矛盾する。

片側対称性を得た後は系\ref{cor:all-defective-one-sym-finite}を適用する。
対称核に階数2以下の非零元があれば中央支持3の点ができるので今の場合にはない。
他の場合に正則なら両対称であるが、定理\ref{thm:all-bisymmetric-as-closure}より両対称な標準正則窓の反対階数は8以下である。
従って中央支持が全て4の候補は除外される。
\end{proof}

\begin{corollary}\label{cor:all-windows-through-eight}
反対階数8以下では、有限Frobenius条件を通る全窓の全次数正則性と、正則な窓どうしの
添字保持同型を有限に判定できる。階数9でこれらの判定の被覆に残る正則候補は、
曲面射影を持たず、純関係と正次元射影ファイバーの双方を持つ窓に限る。
\end{corollary}
\begin{proof}
曲面以上の射影には定理\ref{thm:no-moving-plane-pencils}の被覆を使う。
曲線射影で分岐があれば、階数7以下は定理\ref{thm:low-rank-branch-exclusion}、
階数8は定理\ref{thm:rank-eight-branch-exclusion}で排除される。
階数9で純関係がない場合も第\ref{sec:branch-rank-eleven-floor}節の下界に反する。
点ファイバーで一次元なら系\ref{cor:all-one-dimensional-point-graphs}の全被覆を使う。
四辺が有限なら定理\ref{thm:finite-frobenius-low-opposite-as}、
\ref{thm:finite-rank-eight-all-marks}及び上の定理を用い、同型判定には
定理\ref{thm:finite-all-window-keys}を用いる。
\end{proof}
